\documentclass[11pt,a4paper]{article}

\usepackage[margin=2.3cm]{geometry}
\usepackage{amsmath,amssymb}
\usepackage{graphicx}
\usepackage{booktabs}
\usepackage{array}
\usepackage{xcolor}
\usepackage{enumitem}
\usepackage[hidelinks]{hyperref}
\usepackage{caption}

\emergencystretch=1.5em
\captionsetup{font=small,labelfont=bf,width=.92\textwidth}
\setlength{\parskip}{4pt plus 1pt}
\setlength{\parindent}{0pt}

\newcommand{\code}[1]{\texttt{#1}}
\newcommand{\topk}{\operatorname{TopK}}
\newcommand{\ok}{\checkmark}
\newcommand{\fail}{\texttimes}
\newcommand{\PH}[1]{\textcolor{red}{[#1]}}

\title{
The RBLapSum surrogate in a code-residual stack\\[2pt]
\large Research journal
}
\author{}
\date{Started 2026-09-30 (box \code{lisbon}), following \code{docs/stream-bottleneck-depth-journal.tex}.}

\begin{document}
\maketitle
\vspace{-2em}
\tableofcontents

% ============================================================
\section{The question}
% ============================================================

The code residual (\code{code\_residual=true}) made a 24-layer stream-bottleneck model train with plain hard
\(\topk\) (val CE 1.266 at 20k steps, \(K=512\)).  The next step was to repeat the \(K\)/\(J\) comparison of
\code{docs/kj-vs-hard-topk.tex} -- RBLapSum \(\topk\) with \(J\) extra backward-only candidates against hard
Top-\(K'\) -- in this architecture: 24 layers, \(d=1024\), \(N=4096\), MD, \code{bottleneck\_init=standard},
\code{bottleneck\_decoder\_scale=none}, \code{through\_rank\_kappa}, \(b_0=0\), shared projections, no post-norm, no
value shift, \(T\in\{1,2\}\) (base config \code{/workspace/configs\_lisbon/base\_24L\_cr\_kj.yaml}).

A 60-step smoke run (\(K=32\), \(J=480\), \(T=1\)) trained at 1.10\,s/step but with a gradient norm of
12\,000--15\,000 at every step, against \(\sim\)1.5 for the hard code residual, and a much slower loss (10.96 at
step 60 against 7.28).  The campaign was held back to find the cause.

Notation.  At gate \(\ell\) the input is \(u=c_\ell+\alpha W_E\Delta_\ell\) (carried code plus encoded update), the
scores are \(s_i=|u_i|\), the rank boundary is \(b=s_{(K+1)}\), the pool is the Top-\((K+J)\), \(m_i\) the hard
mask.  RBLapSum's backward is the hard path plus a support term,
\begin{equation}
  \frac{\partial\mathcal L}{\partial u_i}=m_i g_i+\operatorname{sign}(u_i)\bigl(a_i-q_i\textstyle\sum_j a_j\bigr),
  \qquad a_i=g_i\,u_i\,\kappa_i,\quad \kappa_i=\frac{e^{-|s_i-b|/T}}{2T},\quad q_i=\frac{\kappa_i}{\sum_j\kappa_j},
  \label{eq:rbk}
\end{equation}
with \(g=\partial\mathcal L/\partial c_{\ell+1}\) the gradient arriving at the gate's output.

% ============================================================
\section{Entry 1 --- the backward grows toward the input}
% ============================================================

\code{analysis/depth\_probe.py} at initialization, per-block factor of the fc1 weight gradient
\((g_{23}/g_0)^{1/23}\) (below 1: the gradient grows toward the input):

\begin{center}\small
\begin{tabular}{lcc}
\toprule
cell & \(T=1\) & \(T=2\) \\
\midrule
hard \(K=32\) / hard \(K=512\) & 1.020 / 1.011 & \\
RBLapSum \(K=32,J=32\)         & 0.620 & \\
RBLapSum \(K=32,J=480\)        & 0.619 & 0.677 \\
RBLapSum \(K=256,J=256\)       & 0.662 & 0.712 \\
\(K=32,J=224\): \(T=4\), 8, 16 & \multicolumn{2}{c}{0.734, 0.789, 0.849} \\
\(K=32,J=224\), \(T=1\), \code{rblapsum\_support\_scale} 0.3 / 0.1 & 0.908 / 0.992 & \\
\(K=256,J=256\), \(T=1\), \code{rblapsum\_support\_scale} 0.3 / 0.1 & 0.926 / 0.991 & \\
\(K=32,J=224\), \(T=1\), \code{detach} / \code{project} & 0.617 / 0.642 & \\
\(K=32,J=224\), \(T=1\), \code{lapsum} (mass-constrained) & 0.474 & \\
\(K=32,J=224\), \(T=1\), update scale \(\alpha=0.3\) & 0.777 & \\
\bottomrule
\end{tabular}
\end{center}

\textbf{Measured:} with the surrogate the gradient grows by a factor 1.4--1.6 per block toward the input (first
block \(5.7\) against \(10^{-4}\) for the last at \(K=32,J=224,T=1\)).  It depends little on \(J\) and on the
boundary mode, falls only slowly with \(T\) (still 1.18 per block at \(T=16\)), and disappears when the whole
support term is scaled by 0.1.  In the post-norm architecture the same surrogate was near-critical (\(\times1.07\)
per block at \(T=1\), depth journal Entry 3).

% ============================================================
\section{Entry 2 --- where the growth comes from}
% ============================================================

\code{analysis/surrogate\_probe.py} splits the gradient at every gate into its hard part \(m\odot g\) and the
surrogate part (the rest), and records the boundary and the code scale.  At \(K=32\), \(J=224\), \(T=1\):

\begin{center}\small
\begin{tabular}{lcccccc}
\toprule
gate & 0 (entry) & 1 & 6 & 12 & 18 & 24 \\
\midrule
boundary \(b\)                        & 2.67 & 3.63 & 4.15 & 4.25 & 4.33 & 4.39 \\
RMS of the kept scores                & 3.00 & 4.19 & 6.33 & 7.71 & 8.76 & 9.65 \\
pool members with \(|s-b|<T\)          & 255 & 189 & 84 & 63 & 53 & 49 \\
\(\sigma\) = surrogate / hard energy   & 3.77 & 3.47 & 1.79 & 1.34 & 1.20 & 3.11 \\
RMS of \(\partial\mathcal L/\partial c\) (hard model) & \(2.3\cdot10^{-6}\) & \(1.9\cdot10^{-6}\) & \(1.6\cdot10^{-6}\) & \(1.5\cdot10^{-6}\) & \(1.5\cdot10^{-6}\) & \(9.1\cdot10^{-6}\) \\
RMS of \(\partial\mathcal L/\partial c\) (RBLapSum) & 0.60 & 0.24 & \(1.1\cdot10^{-2}\) & \(5.5\cdot10^{-4}\) & \(3.1\cdot10^{-5}\) & \(9.1\cdot10^{-6}\) \\
\bottomrule
\end{tabular}
\end{center}

Three facts follow.  (i) The gradient on the carried code is dense: the next block reads the code through the
decoder, \(x=g_DW_Dc\), so \(\partial\mathcal L/\partial c=g_DW_D^\top\partial\mathcal L/\partial x\) reaches every
coordinate, active or not (in the hard model 90\,\% of its energy is off the support at gate 12).  The hard gate
passes only the \(K\) active coordinates; the surrogate also passes the pool.  (ii) The kept scores are large
compared with \(T\): the boundary sits at 3.6--4.4 and the kept values at 4--10, because the code accumulates
additively like a residual stream.  (iii) The surrogate's energy is comparable to or larger than the hard part at
every gate (\(\sigma=1.2\)--3.8).

\textbf{The mechanism.}  Per coordinate, the backward of one gate multiplies \(g_i\) by
\begin{equation}
  \text{hard: }m_i\in\{0,1\},\qquad
  \text{RBLapSum: }m_i+|u_i|\,\kappa_i\ \ (\text{minus the rank-one correction}),
\end{equation}
and the carry \(u=c_\ell+\dots\) is the identity on the same coordinate.  An active coordinate near the boundary
is therefore multiplied by \(1+|u_i|\kappa_i\approx1+b/(2T)\approx3\) at \(T=1\), at every gate it passes; since
the code carries features from gate to gate, the same coordinates stay near the boundary over several gates and
the factors compound.  In the post-norm architecture every layer re-encodes into fresh coordinates, so no
per-coordinate product forms (and the selection gain's \(1/s\) dominated).  The table's trend in \(T\) fits this:
larger \(T\) lowers \(\kappa\) at the boundary but spreads it over more of the pool, and \(1+|u|\kappa>1\) holds for
every active coordinate at any \(T\).

\textbf{A smoothed forward is not the way out} (\code{rblapsum\_sf}, the soft forward whose backward is its exact
gradient): \(\times0.717\) per block at \(K=32,J=224,T=1\), \(\sigma=4.8\).  The soft gate \(u\mapsto u\,p((|u|-b)/T)\)
has slope \(p+|u|\kappa\), which exceeds 1 in a band around the boundary (about \(1/2+b/(2T)\approx2.6\) at it): the
smoothed stack genuinely amplifies perturbations along the carried coordinates.  The growth is not an artefact of
the mismatch between a hard forward and a soft backward; it comes from any smooth selection with slope above 1,
compounded by an identity carry.

% ============================================================
\section{Entry 3 --- keeping the carry's backward exact}
% ============================================================

What must be preserved is the exact Jacobian of the carry (the mask, isometric on the support), while the
surrogate keeps doing its job: giving gradient to the \(J\) candidates so the model learns what to select.  Two
ways, both prototyped as patches (\code{surr\_fix\_proto.py}, \code{train\_patched.py} on the box):

\begin{itemize}[leftmargin=1.2em]
\item \textbf{(F1) local surrogate.}  The support term is routed only into the block's own update
  \(\alpha W_E\Delta_\ell\); the carried code receives the exact hard gradient \(m\odot g\).  Each selection
  decision is credited once, to the block that makes it, rather than once per gate the feature passes.
  Implemented as \(y=\mathrm{gate}(\bar c+v)+\bigl(y_h-\bar y_h\bigr)\) with \(\bar{\cdot}\) a detach and
  \(y_h=(c+\bar v)\odot m\): the forward is unchanged, the carry gets \(m\odot g\), the update gets the full
  surrogate gradient.
\item \textbf{(F3) inactive-only surrogate.}  The support term is kept for the inactive pool members and dropped on
  active ones, which then keep the exact multiplier 1 (the \(\kappa\)-weighted correction is taken over the
  inactive members).
\end{itemize}

Per-block factor of \(\partial\mathcal L/\partial c\) at initialization (hard: \(\times1.07\), unmodified: \(\times0.64\)):

\begin{center}\small
\begin{tabular}{lccc}
\toprule
 & \(K=32,J=224,T=1\) & \(K=32,J=224,T=2\) & \(K=256,J=256,T=1\) \\
\midrule
unmodified     & 0.643 (\(\sigma=1.72\)) & 0.704 (0.66) & 0.683 (0.87) \\
(F1) local     & 0.999 (\(\sigma=1.26\)) & 1.020 (0.57) & 1.000 (0.51) \\
(F3) inactive  & 1.053 (\(\sigma=0.37\)) & 1.064 (0.17) & 1.043 (0.09) \\
\bottomrule
\end{tabular}
\end{center}

\textbf{Measured:} both remove the growth.  F1 keeps the surrogate's full energy at the gate (it is only routed
away from the carry); F3 keeps a third to a tenth of it.  \textbf{Interpretation:} the diagnosis holds --- the
growth is the surrogate term compounding along the carried coordinates.

% ============================================================
\section{Entry 4 --- training screens}
% ============================================================

2000-step screens at \(K=32\), \(J=224\) unless marked, 24 layers, the campaign's base config; the fixes were
applied with the prototype entry \code{train\_patched.py} (identical to \code{rblapsum\_surrogate\_scope}, landed
later in \code{4beaf1b}).  Hard code residual at \(K=32\) (\code{li\_24L\_k32\_coderes\_a1}) for reference.

\begin{center}\small
\begin{tabular}{lcccccc}
\toprule
 & \multicolumn{3}{c}{val CE at step} & & & \\
run & 500 & 1000 & 1500 & dead features & usage entropy & boundary \(b\) \\
\midrule
hard \(K=32\) (reference)          & 2.912 & 2.453 & 2.241 & 0.001 & 0.41 & -- \\
unmodified, \(T=1\)                & 5.784 & 5.496 & 5.340 & 0.94 & 0.035 & 105--335 \\
\code{rblapsum\_support\_scale} 0.1 & \textbf{2.890} & \textbf{2.428} & \textbf{2.216} & 0.000 & 0.64 & 3.9--2.4 \\
(F3) inactive only, \(T=1\)         & 2.912 & 2.467 & 2.256 & 0.006 & 0.45 & 5.3--4.3 \\
(F1) local, \(T=1\)                 & 4.908 & 5.443 & 5.629 & 0.93 & 0.031 & 93--799 \\
(F1) local, \(T=2\)                 & 3.467 & 5.208 & 5.344 & & & \\
(F1) local, \(J=32\), \(T=1\)        & 4.072 & 5.395 & 5.194 & & & \\
\bottomrule
\end{tabular}
\end{center}

\textbf{Measured.}  The unmodified surrogate and F1 both collapse onto \(\sim\)6\,\% of the features (usage entropy
0.03), whose values run away (the rank boundary grows from 4 to hundreds), and the loss rises.  F1 does so although
its backward was flat at initialization.  F3 is stable and equals hard within 0.015.  A global support scale of 0.1
is stable, spreads the usage (entropy 0.64 against 0.41 for hard) and is ahead of hard by 0.02--0.025 at every
evaluation.

\textbf{What F1 taught.}  Removing the compounding from the backward is not enough: the support term
\(a_i=g_i\,u_i\,\kappa_i\) is proportional to the value \(u_i\).  For an \emph{active} feature it makes the gradient on
its value proportional to the value itself --- a growth-rate term, \(\dot u\propto u\) --- and in a code that
accumulates through an identity carry, that is exponential growth of whichever features the loss favours, until
they win every selection.  Routed into the update (F1) or into the carry (unmodified), the term runs away the same
way.  The post-norm architecture never showed this because every layer reset the scale.  With an absolute \(T\)
there is a second feedback: as the code grows, \(\kappa\) at the boundary stays \(1/(2T)\) while \(u\approx b\) grows, so
the push itself grows with the code scale.  F3 removes the active-feature term altogether and is stable, but gains
nothing; a small dose of it on every member (scale 0.1) is stable and helps a little.

% ============================================================
\section{Entry 5 --- the code's scale sets the surrogate's strength}
% ============================================================

\textbf{Measured} (the logs of Entry 4; \(b\) is the rank boundary averaged over the 25 gates and the batch, the
usage entropy is normalized to 1 for uniform use):

\begin{center}\small
\begin{tabular}{lcccccccc}
\toprule
 & \multicolumn{8}{c}{\(b\) \;/\; usage entropy at step} \\
run & 1 & 100 & 200 & 300 & 400 & 500 & 1000 & 2000 \\
\midrule
(F1) local, \(T=1\) & 4.16/0.80 & 5.48/0.44 & 7.14/0.17 & 13.2/0.09 & 40.7/0.06 & 93.4/0.04 & 388/0.03 & 1179/0.03 \\
(F1) local, \(T=2\) & 4.16/0.80 & 5.66/0.46 & 6.56/0.30 & 8.85/0.15 & 14.2/0.14 & 44.2/0.11 & 331/0.04 & 1911/0.03 \\
unmodified, \(T=1\) & 4.16/0.80 & 7.29/0.39 & 14.1/0.06 & 37.1/0.04 & 64.3/0.04 & 105/0.03 & 189/0.04 & 470/0.04 \\
support scale 0.1   & 4.16/0.80 & 5.15/0.48 & 4.53/0.62 & 4.32/0.67 & 4.15/0.69 & 3.86/0.68 & 2.76/0.65 & 2.23/0.62 \\
(F3) inactive only  & 4.16/0.80 & 5.71/0.42 & 5.09/0.51 & 5.00/0.55 & 5.16/0.56 & 5.26/0.54 & 4.65/0.49 & 4.27/0.42 \\
\bottomrule
\end{tabular}
\end{center}

All runs lose usage entropy in the first 100 steps and their boundary rises to 5--6.  In the stable runs the
boundary then falls and the entropy partly recovers; in F1 and the unmodified run the boundary keeps growing,
accelerating (F1 at \(T=1\): \(\times1.3\), \(\times1.9\), \(\times3.1\), \(\times2.3\) per 100 steps from step 100), while the
usage concentrates and features die.  F1's training CE follows the stable runs to step \(\sim\)150 (6.48 at 100
against 6.42--6.41) and departs as the boundary grows (4.13 at 200 against 4.07--4.08; 3.70 at 300 against
3.44--3.45; rising from 400).

\textbf{Interpretation (a hypothesis to test).}  Hard \(\topk\) is 1-homogeneous, \(\topk(\lambda u)=\lambda\topk(u)\).
Every block reads the code through its pre-LN RMSNorm and the final norm removes the code's scale, so the model's
function depends on the code scale only through the ratio of each update to the carry; the scale is weakly
constrained.  The surrogate with a fixed \(T\) is not scale-free: at the boundary \(|u|\kappa=b/(2T)\), and the exact
surrogate gain of \code{docs/scale-dynamics-note-neutral.tex},
\begin{equation}
  \Pi=\bigl\lVert L_\kappa D_z\bigr\rVert_F^2\approx\frac{b^2}{4T\delta},\qquad
  L_\kappa=D_\kappa-\frac{\kappa\kappa^\top}{\sum_i\kappa_i},
\end{equation}
(\(\delta\) the local rank spacing at the boundary) grows linearly with the scale when the shape of the score
distribution is fixed (\(b\) and \(\delta\) both scale).  That note found that a large \(\Pi\) drives large support
updates that raise the score scale at second order in the step (``heating'').  With \(T\) fixed, a larger scale is a
larger \(\Pi\), and the loop closes.  In the post-norm stack each gate's input is normalized; the code residual has
no such normalization.

\textbf{The fix to test: a scale-free kernel.}  \code{rblapsum\_relative\_temperature} (commit \code{2914577}) sets
the kernel width to \(T_{\rm eff}=\tau\,b\) per row.  Then \(|u|\kappa=1/(2\tau)\) at the boundary at every scale,
\(\Pi\approx b/(4\tau\delta)\) depends on the shape of the scores only, and scaling the gate input is an exact
reparameterization of the whole backward (unit test).  At initialization \(b\approx4.2\), so \(\tau=0.25\) is the
\(T=1\) kernel: \textbf{prediction} --- if the loop drives the collapse, F1 at \(\tau=0.25\) does not collapse, although
it starts from exactly the state of the F1 run that did.  \code{tools/repro\_check.py}: IDENTICAL against
\code{4beaf1b} (hard and rbk), so the default path is unchanged.

Round 2 (2000 steps, \(K=32\), \(J=224\), same base config):

\begin{center}\small
\begin{tabular}{lll}
\toprule
run & change & tests \\
\midrule
\code{upd\_rel025} & F1, \(\tau=0.25\) & the loop (same kernel as F1 \(T=1\) at step 0) \\
\code{upd\_rel05}, \code{upd\_rel1} & F1, \(\tau=0.5\), 1 & strength at fixed shape \\
\code{upd\_t1\_ss03} & F1, \(T=1\), support scale 0.3 & strength alone, without scale-freedom \\
\code{pool\_t1\_ss03}, \code{pool\_t1\_ss003} & unmodified, support scale 0.3, 0.03 & the existing knob (0.1 was best so far) \\
\code{pool\_t1\_ss0} & unmodified, support scale 0 & control: hard backward, boundary logged \\
\bottomrule
\end{tabular}
\end{center}

\textbf{Measured} (val CE at steps 500/1000/1500/2000; usage entropy and dead fraction at step 600; \(\sigma\): surrogate/hard
energy at the gate, mean over gates, \code{surrogate\_probe.py} at initialization):

\begin{center}\small
\begin{tabular}{lcccc}
\toprule
run & \(\sigma\) at step 0 & val CE 500 / 1000 / 1500 / 2000 & entropy / dead (600) & \(b\) (600) \\
\midrule
support scale 0 (exact hard backward)     & 0     & 2.896 / 2.435 / 2.232 / 2.104 & 0.61 / 0.000 & 4.2 \\
F1, \(\tau=0.25\)                         & 1.37  & 2.948 / 2.675 / -- / --       & 0.12 / 0.48 & 14.9 \\
F1, \(\tau=0.5\)                          & 0.61  & 2.922 / 2.549 / -- / --       & 0.15 / 0.40 & 11.9 \\
F1, \(\tau=1\)                            & 0.26  & 2.899 / 2.491 / 2.350 / 2.258 & 0.18 / 0.30 & 10.7 \\
F1, \(T=1\), support scale 0.3            & 0.087 & 2.871 / 2.424 / 2.250 / 2.189 & 0.26 / 0.008 & 5.5 \\
unmodified, support scale 0.3            & 0.11  & 2.937 / 2.440 / 2.206 / \textbf{2.064} & 0.37 / 0.001 & 4.6 \\
unmodified, support scale 0.1 (Entry 4)  & 0.010 & 2.890 / 2.428 / 2.216 / 2.089 & 0.67 / 0.000 & 4.3 \\
unmodified, support scale 0.03           & 0.001 & 2.893 / 2.436 / 2.226 / 2.094 & 0.63 / 0.000 & 3.7 \\
F3, \(T=1\) (Entry 4)                    & 0.37  & 2.912 / 2.467 / 2.256 / 2.119 & 0.53 / 0.002 & 5.2 \\
\bottomrule
\end{tabular}
\end{center}

\textbf{The prediction failed.}  F1 with the scale-free kernel collapses at every width: \(\tau=0.25\), which starts
from the same kernel as the F1 \(T=1\) run, loses features at nearly the same pace (dead 0.40 at step 400 against
0.76), and \(\tau=1\), with a quarter of the gain at the boundary (\(|u|\kappa=0.5\)) and \(\sigma=0.26\), still loses 30\,\%
of the features by step 600.  The growing scale is not what drives the collapse; with the scale removed from the
kernel the boundary still grows (4.2 to 11--15).  The width matters little and the total strength a lot: F1 at
\(T=1\) with the whole support term scaled by 0.3 (\(\sigma=0.087\)) is the best run at step 1000, though its usage
entropy also falls.  The F1 \(\tau=0.25\) and 0.5 runs were stopped at step 1140 (collapsing).  F1 at support scale 0.3 leads
at step 1000 and then falls behind (2.189 at 2000, +0.085), its usage still concentrating.  Without the update
routing, the unmodified surrogate at support scale 0.3 is the best run at step 2000, 0.040 below the exact hard
backward, ahead of 0.1 (\(-0.014\)) and 0.03 (\(-0.009\)); its usage also concentrates (entropy 0.19 and 4\,\% dead
features at step 2000), but at step 2000 it holds at most 10 always-on features (Entry 6) and \(f\le0.12\).  This
run and support scale 0.1 continue to 20k steps.

\textbf{Correction to Entry 4.}  The hard reference of Entry 4 used \code{bottleneck\_decoder\_scale=backward\_preserving}
(the depth study's base); the campaign base uses \code{none}.  The exact hard backward on the campaign base
(support scale 0) reaches 2.896 at step 500 against 2.912: the 0.02 advantage of support scale 0.1 in Entry 4 is
mostly this difference.  The support-scale-0 run is the reference from here on.

% ============================================================
\section{Entry 6 --- what fills the support: always-on features}
% ============================================================

\code{code\_usage.py} (on the box; 4 validation batches of \(8\times512\) tokens, no gradient) records the code
\(c\) after every gate: \(n_{>1/2}\), the number of features active for more than half of the tokens; \(E_{>1/2}\),
their share of the code energy; \(f\), the common-mode fraction of the code within a sequence (as in the depth
study, positions \(t\ge16\)); and the rank boundary \(b\).  Gates 6 / 12 / 18 / 24:

\begin{center}\small
\begin{tabular}{lcccc}
\toprule
checkpoint & \(n_{>1/2}\) & \(E_{>1/2}\) & \(f\) & \(b\) \\
\midrule
F1, \(T=1\), step 2000 (CE 5.9)        & 32 / 32 / 32 / 32 & .95 / .93 / .91 / .95 & .84 / .83 / .81 / .86 & 620 / 1015 / 3622 / 6145 \\
unmodified, \(T=1\), step 2000 (5.3)   & 27 / 25 / 26 / 24 & .67 / .73 / .77 / .79 & .53 / .67 / .71 / .80 & 175 / 520 / 149 / 1935 \\
support scale 0.1, step 2000           & 0 / 0 / 0 / 0     & 0 / 0 / 0 / 0         & .01 / .01 / .01 / .05 & 1.7 / 1.9 / 2.1 / 7.9 \\
F3 (inactive only), step 2000          & 0 / 1 / 2 / 10    & 0 / .02 / .10 / .29   & .01 / .01 / .08 / .18 & 2.5 / 2.7 / 3.8 / 17 \\
hard \(K=32\), step 16k                & 0 / 1 / 1 / 2     & 0 / .01 / .01 / .05   & .02 / .02 / .02 / .06 & 1.8 / 2.1 / 2.5 / 6.7 \\
hard \(K=512\), step 20k               & 30 / 59 / 88 / 83 & .04 / .08 / .14 / .17 & .04 / .04 / .06 / .07 & 1.6 / 2.0 / 2.3 / 3.6 \\
\bottomrule
\end{tabular}
\end{center}

\textbf{Measured.}  In the collapsed F1 run every one of the \(K=32\) slots from gate 6 on is held by a feature that
is active for more than half of the tokens; they carry 91--95\,\% of the code's energy, the codes of all positions
nearly coincide (\(f=0.81\)--0.86), and their values reach \(10^2\)--\(10^4\), growing with depth.  The unmodified run
is the same with 24--27 such features.  The hard models keep at most 2 of 32 (and 17\,\% of the energy at
\(K=512\)); support scale 0.1 has none; F3 has 10 at the last gate holding 29\,\% of its energy --- the same failure,
slower.

\textbf{Interpretation: a token collapse in code space.}  The code is the residual stream, and it stops depending on
the token when the \(K\) slots are taken by features active for every token; the readout then sees only their
token-independent values (CE \(\approx\) the unigram entropy, 5.97).  The surrogate creates them because it gives an
\emph{inactive} feature a gradient on its score, with a sign set by how much the loss would gain from including it.
A generic feature --- a direction useful for every token, which early in training is the unigram --- receives a push
of the same sign at every token, and the parameters that implement a token-independent score shift (the encoder row
against the mean of the updates) raise it for all tokens at once, until it is active everywhere.  In the code
residual every block then adds its token-common update to an always-on feature, so its value grows with depth and
with training, far above the kernel window, where neither the surrogate nor bounded optimizer steps bring it back.
The hard gradient gives an inactive feature no score gradient at all: how often a feature is used follows only from
value learning, and the hard models end with few always-on features holding little energy.

\textbf{Fix to test (F4): center the support term over tokens.}  The boundary correction of
\code{through\_rank\_kappa} already removes one common mode of the support term: the mean over one token's pool, a
shift of all of that token's scores, which only moves its boundary.  The collapse runs through the other one: the
mean over the tokens of one feature's term, a shift of that feature's score for all tokens, which changes how often
it is used.  \code{rblapsum\_center\_tokens} (commit \code{ced6be9}) removes it: per gate and feature, the mean of the
support term over the micro-batch's tokens whose pool contains the feature is subtracted on those tokens.  The
surrogate then decides \emph{which} tokens use a feature; \emph{how often} it is used is left to the exact gradient,
as in the hard model.  Oracle: IDENTICAL against \code{4beaf1b}.  \textbf{Prediction:} F1 and F3 at full strength
with centering keep \(n_{>1/2}\) near the hard model's and do not collapse.

Round 3 (2000 steps, \(K=32\), \(J=224\), \(T=1\)):

\begin{center}\small
\begin{tabular}{lll}
\toprule
run & change & tests \\
\midrule
\code{updinact\_t1} & F1 routing, inactive members only & which half of the term drives the collapse \\
\code{updact\_t1} & F1 routing, active members only & \\
\code{upd\_t1\_ctr} & F1 + token centering & the fix at full strength \\
\code{inact\_t1\_ctr} & F3 + token centering & \\
onset runs & F1, support scale 0, F1 + centering; a checkpoint every 50 steps to 600 & the timeline of \(n_{>1/2}\), \(f\), \(b\) \\
\bottomrule
\end{tabular}
\end{center}

\textbf{Measured: the onset} (F1, \(T=1\), a checkpoint every 50 steps; \(n_{>1/2}\) / \(E_{>1/2}\) / \(f\)):

\begin{center}\small
\begin{tabular}{lcccc}
\toprule
step (train CE) & gate 6 & gate 12 & gate 18 & gate 24 \\
\midrule
50 (9.03)   & 0 / .00 / .05  & 0 / .00 / .06  & 5 / .15 / .16  & 18 / .65 / .62 \\
100 (5.83)  & 0 / .00 / .04  & 0 / .00 / .04  & 3 / .07 / .08  & 24 / .70 / .65 \\
200 (3.89)  & 25 / .48 / .46 & 8 / .06 / .08  & 9 / .10 / .09  & 18 / .31 / .24 \\
300 (3.47)  & 31 / .87 / .85 & 15 / .11 / .12 & 21 / .27 / .26 & 14 / .22 / .20 \\
450 (3.92)  & 32 / .96 / .94 & 32 / .83 / .81 & 32 / .92 / .91 & 12 / .23 / .24 \\
\bottomrule
\end{tabular}
\end{center}

The takeover starts at the \emph{last} gate within 50 steps, while the model is learning the unigram (train CE 9),
reaches gate 6 by step 200 and the middle gates by step 450; the last gate partly recovers.

The same series for the exact hard backward (support scale 0): the last gate forms always-on features in the
unigram phase as well (10 features with 60\,\% of the energy at step 50, 18 with 69\,\% at step 100, \(f=0.64\)), and
they dissolve as training proceeds (6 with 10\,\% at step 600); gates 6, 12 and 18 never have any (\(f\le0.05\)).
\textbf{Interpretation:} an always-on code at the readout during the unigram phase is the ordinary way the model
represents the unigram.  What the surrogate adds is the spread of always-on features to the lower gates and their
lock-in.

\textbf{Measured: round 3} (val CE at steps 500 / 1000 / 2000; the support-scale-0 run for reference: 2.896 / 2.435 /
2.104):

\begin{center}\small
\begin{tabular}{lccc}
\toprule
run & val CE & dead (step 1000) & note \\
\midrule
F1 routing, active members only   & 3.141 / 3.025 / 4.918 & 0.73 & collapse \\
F1 routing, inactive members only & 2.920 / 2.525 / 2.280 & 0.01 & behind by 0.18 at 2k \\
F1 + token centering              & 5.263 at 500          & 0.93 at step 200 & stopped at step 500 \\
F3 + token centering              & 5.778 at 500          & 0.90 at step 200 & stopped at step 500 \\
\bottomrule
\end{tabular}
\end{center}

\textbf{The ablation.}  Both halves of the routed term hurt; the active half is the destructive one (73\,\% of the
features dead by step 1000, against 1\,\%; val CE 4.918 at step 2000), and the inactive half alone is 0.18 behind
the hard backward at step 2000, with the boundary growing steadily (5.0 to 9.5).  \textbf{Interpretation:} for an
active feature the routed term multiplies the update's gradient by \(1+|u|\kappa\approx3\) at the boundary while the
carry keeps \(m\odot g\); block \(\ell\) then learns the values of its marginal features three times faster than the
blocks that wrote them, see Entry 7.

\textbf{The prediction for F4 failed, badly.}  Token centering does not prevent the collapse; it accelerates it: dead
fraction 0.11 at step 100 and 0.93 at step 200, the boundary 4 to 12 by step 100, gradient norm \(\sim\)10 from step 40
(the uncentered F1 reached these states only at step 300--500).  The per-token zero-sum survives centering (the
logged \(|\sum_i g_{s,i}|/\lVert g_s\rVert\) stays \(\le10^{-4}\)), and at initialization centering changes \(\sigma\) by
less than 1\,\% (the token mean of the support term is 0.4\,\% of its energy, the level of chance for \(\sim\)256
tokens per feature).  \textbf{Measured} (the centered F1 run with a checkpoint every 50 steps, \(n_{>1/2}\) / \(E_{>1/2}\) / \(f\) at gates
6 / 12 / 18 / 24): at step 50 \(f=0.12\) / 0.13 / 0.22 / 0.44 (uncentered F1: 0.05 / 0.06 / 0.16 / 0.62); at step 100
every gate holds 30--31 always-on features with 91--98\,\% of the energy and \(f=0.90\)--0.96, at all gates at once,
where the uncentered run started at the last gate and needed 450 steps to reach the middle ones.  On checkpoints of
the uncentered runs the support term's per-feature token mean is at the level of chance (0.4--0.6\,\% of its energy
at gates 0--18, 2--5\,\% at the last gate) and uncorrelated with how often a feature is used (\(|r|\le0.09\)).

\textbf{Interpretation.}  Centering is not a small change even though the mean it removes is small: it moves that
mean from the few tokens where the feature sits near the boundary to every token whose pool contains the feature,
including those where it is active far above the boundary and the kernel is zero.  There the term acts on the
feature's value, with the same sign at every token, so the gradient on each block's output acquires a
token-independent component, \(W_E^\top(\text{per-feature constant})\), at every gate at once; a token-independent
update is exactly what feeds always-on features.  The support term must stay where the kernel puts it.
Token centering is abandoned.

% ============================================================
\section{Entry 7 --- one support term per gradient path}
% ============================================================

Two routings have now been tested against the compounding of Entries 1--3: the unmodified one (each gate's support
term applied to the total upstream, which already contains the terms of later gates, so they multiply along the
carry) and F1 (the term routed into the update only: no compounding, but a block gets no credit for how the
features it writes are selected further up).  F1 trains worse than the unmodified surrogate at every support scale
tried (0.3: 2.189 against 2.064 at step 2000).  The estimator the stack actually calls for sits between them.  A
hard \(\topk\) stack is piecewise linear; the surrogate estimates what a change of support would do.  To first
order in the support changes, one change at a time, the gradient is
\begin{equation}
  \frac{\partial\mathcal L}{\partial\theta}\;\approx\;\underbrace{J_{\rm hard}^\top g}_{\text{no support change}}
  \;+\;\sum_\ell \Bigl(\frac{\partial u_\ell}{\partial\theta}\Bigr)^{\!\top}_{\!\rm hard} S_\ell^\top H_\ell ,
\end{equation}
where \(H_\ell\) is the gradient that reaches gate \(\ell\)'s output along hard paths only, \(S_\ell^\top\) the gate's
support term, and every other Jacobian is the hard one.  The carry and the update receive gate \(\ell\)'s term alike
(unlike F1), and it is never multiplied by another gate's (unlike the unmodified backward, whose extra terms are
products of support changes at several gates, second and higher order).

Stated as two requirements on any surrogate backward for this stack: \textbf{(R1) one gradient for the gate's
input} --- \(c_\ell\) and \(v_\ell\) enter the gate only through \(u=c_\ell+v_\ell\), so an estimator of
\(\partial\mathcal L/\partial u\) must give both the same gradient; F1 violates it (the carry gets \(m\odot g\), the update
\(m\odot g+S^\top g\)), so block \(\ell\) and the blocks that wrote the carried code are trained on different estimates
of the same derivative.  \textbf{(R2) no products of support terms} --- the hard forward has no such terms; the
unmodified backward violates it, and the products grow along the carry.  The unmodified backward at support scale
\(s\) satisfies R1 and violates R2 only at order \(s^2\), which is the reading of why \(s=0.1\)--0.3 trains; F3 satisfies
R1 and keeps products only along stretches where a feature stays just below the boundary
(\(|u|\kappa<1\) there).  The first-order estimator satisfies both at full strength.

\code{rblapsum\_surrogate\_scope=first\_order} (commit \code{2965fa6}) computes it with two backward passes per
micro-batch: a hard pass (every gate passes back \(m\odot g\) and keeps \(g=H_\ell\)), as
\code{torch.autograd.grad} to the token embedding so that no parameter's \texttt{.grad} is touched, then the
ordinary pass, in which each gate adds the support term computed from its kept \(H_\ell\).  Unit tests: against the
explicit two-gate formula, and, on a 3-layer code-residual model (with its read paths), against the hard gradient
plus the sum over the gates of (gradient with only that gate relaxed \(-\) hard gradient), equal to float64
round-off; oracle IDENTICAL against \code{4beaf1b} (hard and rbk).  The extra pass costs about 45\,\%
of a step (1.57 against 1.09\,s at \(K=32\)).  At initialization the per-block factor
of \(\partial\mathcal L/\partial c\) is \(\times0.982\) (hard 1.072, unmodified 0.643): what remains is the additive
sum of the later gates' terms, not a product.

Round 4: first order at support scale 1 and 0.3 (\(K=32\), \(J=224\), \(T=1\)), 2000 steps.

\begin{center}\small
\begin{tabular}{lccccc}
\toprule
run & val CE 500 / 1000 / 1500 / 2000 & \(-\) hard at 2k & dead (2k) & entropy (2k) & \(b\) (2k) \\
\midrule
hard backward (support scale 0)   & 2.896 / 2.435 / 2.232 / 2.104 & 0 & 0.000 & 0.56 & 2.7 \\
\textbf{first order, scale 1}     & 2.912 / 2.417 / 2.174 / \textbf{2.028} & \(\mathbf{-0.076}\) & 0.000 & 0.36 & 2.1 \\
first order, scale 0.3            & 2.877 / 2.410 / 2.195 / 2.054 & \(-0.050\) & 0.000 & 0.48 & 2.1 \\
unmodified, scale 0.3             & 2.937 / 2.440 / 2.206 / 2.064 & \(-0.040\) & 0.04 & 0.19 & 2.7 \\
unmodified, scale 0.1             & 2.890 / 2.428 / 2.216 / 2.089 & \(-0.014\) & 0.000 & 0.62 & 2.2 \\
\bottomrule
\end{tabular}
\end{center}

At step 2000 the first-order run has at most 2 always-on features per gate (\(E_{>1/2}\le0.05\), \(f\le0.15\) at the
last gate, \(\le0.05\) below).

\textbf{Measured.}  The first-order estimator trains at full strength, with no dead features, the boundary falling
like the hard model's, and is the best run at steps 1000, 1500 and 2000: 0.076 below the exact hard backward at step
2000, against 0.040 for the best rescaling of the unmodified backward.  At the same scale 0.3 the first-order
estimator is ahead of the unmodified one by 0.010.  Its feature usage concentrates (entropy 0.36 against 0.56) without
dead features or always-on features.

\textbf{Interpretation.}  The estimator the code-residual stack needs is the sum of single-gate relaxations, not
their product: the product's extra terms (R2) were what made the unmodified surrogate fail, and the routing's
two gradients for one input (R1) what made the routed one fail.  The first-order estimator is the hard gradient plus the sum over the gates of
(the backward pass in which only gate \(\ell\) uses its surrogate Jacobian \(-\) the hard one); every one of these is a
backward pass through the gate's input \(u_\ell\), so R1 holds by construction.  Damping the product (support scale) approximates it to first order in the
scale and pays for the second-order terms with a smaller first-order term.

Continuations: the first-order run (scale 1) to 20k steps; the hard backward and the unmodified backward at scales
0.1 and 0.3 are already continuing.  Val CE minus the hard backward's, at steps 1k / 2k / 3k / 4k / 5k / 6k / 7k / 8k / 9k:

\begin{center}\small
\begin{tabular}{lccccccccc}
\toprule
 & 1k & 2k & 3k & 4k & 5k & 6k & 7k & 8k & 9k \\
\midrule
first order, scale 1   & \(-0.018\) & \(-0.076\) & \(-0.074\) & \(-0.068\) & \(-0.065\) & & & & \\
unmodified, scale 0.3  & \(+0.006\) & \(-0.040\) & \(-0.048\) & \(-0.049\) & \(-0.046\) & \(-0.036\) & \(-0.033\) & \(-0.028\) & \\
unmodified, scale 0.1  & \(-0.006\) & \(-0.014\) & \(-0.018\) & \(-0.015\) & \(-0.015\) & \(-0.013\) & \(-0.011\) & \(-0.009\) & \(-0.008\) \\
\bottomrule
\end{tabular}
\end{center}

\textbf{Measured:} every surrogate's advantage peaks between 2k and 4k steps and then shrinks; the hard backward is
catching up.

\textbf{The 20k values} (val CE, difference to the hard backward): hard 1.5020; first order 1.4569 (\(-0.045\); \(-0.050\)
at 10k, \(-0.048\) at 16k); first order at \(T=2\) 1.4657 (\(-0.036\); \(-0.088\) at 2k); unmodified at scale 0.3 1.4856
(\(-0.016\)); at scale 0.1 1.4961 (\(-0.006\)).  The first-order advantage shrinks during the first 10k steps and is
nearly constant after; the damped unmodified backward keeps little.  A second seed (initialization and data order)
at step 2000: hard 2.0821, first order 2.0254 (\(-0.057\)), first order at \(T=2\) 2.0125 (\(-0.070\)); first seed
\(-0.076\) and \(-0.088\).  Generality screens: first order at \(K=256\), \(J=256\) (with its hard
control), at \(J=480\), and at \(T=2\); first order with per-block projections.

% ============================================================
\section{Entry 8 --- does the first-order estimator carry over?}
% ============================================================

2000-step screens of the first-order estimator (support scale 1) away from \(K=32\), \(J=224\), \(T=1\), each against
the exact hard backward of the same configuration (\(J\) and \(T\) do not enter the hard backward, so the
\(K=32\) control serves all \(K=32\) rows):

\begin{center}\small
\begin{tabular}{lccc}
\toprule
configuration & first order, val CE 500 / 1000 / 1500 / 2000 & hard backward at 2000 & difference \\
\midrule
\(K=32\), \(J=224\), \(T=1\) (Entry 7)      & 2.912 / 2.417 / 2.174 / 2.028 & 2.104 & \(-0.076\) \\
\(K=32\), \(J=224\), \(T=1\), per-block projections & 2.869 / 2.369 / 2.134 / 1.992 & -- & -- \\
\(K=256\), \(J=256\), \(T=1\)               & 2.869 / 2.350 / 2.098 / 1.958 & 1.870 & \(+0.088\) \\
\(K=32\), \(J=480\), \(T=1\)                & 2.919 / 2.430 / 2.184 / 2.040 & 2.104 & \(-0.064\) \\
\(K=32\), \(J=224\), \(T=2\)                & 2.877 / 2.391 / 2.158 / 2.016 & 2.104 & \(-0.088\) \\
\(K=256\), \(J=256\), \(T=1\), inactive members only & 2.798 / 2.257 / 2.015 / 1.891 & 1.870 & \(+0.020\) \\
\(K=256\), \(J=256\), \(T=1\), support scale 0.3 & 2.816 / 2.284 / 2.048 / 1.925 & 1.870 & \(+0.054\) \\
\(K=256\), \(J=256\), \(T=0.3\,b\) (scale-free) & 2.850 / 2.329 / 2.085 / 1.949 & 1.870 & \(+0.078\) \\
\bottomrule
\end{tabular}
\end{center}

\textbf{Measured: \(K=256\).}  At \(K=256\), \(J=256\) the first-order estimator at full strength is \emph{worse} than
the hard backward at every evaluation (\(+0.111\), \(+0.129\), \(+0.109\), \(+0.088\) at steps 500--2000).  Its code holds many
features that are active for more than half of the tokens, 164 / 153 / 67 of the 256 slots at gates 6 / 12 / 18
(hard: 5 / 10 / 15), but they carry little energy (10--15\,\%) and the codes of different positions stay distinct
(\(f\le0.15\)); the rank boundary is lower (1.6 against 2.7--3.2) and the feature-usage entropy falls (0.51 against 0.77).
At initialization its surrogate energy is lower than at \(K=32\) (\(\sigma=0.080\) against 0.18) and the gradient
reaching the code grows \(\times1.08\) per block relative to hard, as at \(K=32\).  On the step-2000 checkpoints the
per-feature token mean of the support term is 0.5--1.4\,\% of its energy and uncorrelated with usage (\(|r|\le0.02\)),
close to chance.

The code's scale at the boundary, step-2000 checkpoints, \code{surrogate\_probe.py} (mean over gates; window =
pool members with \(|s-b|<T=1\)):

\begin{center}\small
\begin{tabular}{lccc}
\toprule
checkpoint & boundary \(b\) & RMS of the kept scores & pool members in the window \\
\midrule
hard backward, \(K=32\)          & 2.71 & 10.8 & 173 of 256 \\
first order, \(K=32\)            & 2.08 & 5.1  & 227 of 256 \\
hard backward, \(K=256\)         & 3.26 & 11.3 & 23 of 512 \\
first order, \(K=256\)           & 2.08 & 4.9  & 428 of 512 \\
\bottomrule
\end{tabular}
\end{center}

With the first-order estimator the code is half as large and its scores crowd around a lower boundary, at both
\(K\); with \(K=256\) the hard model keeps its boundary region nearly empty (23 of 512 within \(\pm T\)), the first-order
model fills it (428).

\textbf{Interpretation (to be tested).}  The always-on features at \(K=256\) are not the token collapse of Entry 6:
they are weak and many, sitting just above a boundary that has moved down.  This is what a compressed boundary
band looks like (\code{docs/rblapsum-kappa-stability-neutral.tex}): the support term pushes marginal active features
down and near candidates up, toward the boundary, and with \(K+J=512\) of 4096 features in the pool many features end
up just above it for most tokens, occupying slots with little content.  With a fixed \(T\), a crowded,
smaller code puts more of the pool inside the kernel, which strengthens the push: a loop.  Tests: the same at support
scale 0.3; the first-order estimator restricted to the inactive members (no push on active features); and the
first-order estimator with the scale-free kernel (\(T=0.3\,b\), the \(T=1\) kernel at initialization).

\textbf{Measured: the other screens.}  At \(J=480\) the first-order estimator gains 0.064 at step 2000 (0.076 at
\(J=224\)), with at most 2 always-on features per gate.  With an encoder/decoder pair per block it reaches 1.992 at
step 2000 (its hard control was dropped for GPU time; the shared-pair comparison is the one above).  At \(K=256\) a
smaller support scale does not rescue it: at scale 0.3 the first-order estimator is \(+0.059\), \(+0.063\), \(+0.059\)
behind the hard backward at steps 500, 1000, 1500, the unmodified backward \(+0.118\), \(+0.128\), \(+0.112\), \(+0.098\)
at steps 500--2000.  The first-order estimator restricted to the inactive members is the closest, \(+0.041\) at step
500, with the usage entropy of the hard model (0.85 at step 400).
\textbf{Measured: how crowded the hard model's boundary is.}  Pool members within \(\pm T\) of the boundary, hard
backward: at \(K=32\) 173 / 193 / 208 of 256 at steps 2k / 4k / 6k (the region fills further during training); at
\(K=256\) 23 of 512 at step 2k (it empties, from 310 at initialization).  The first-order model at \(K=32\): 227 and
235 of 256 at steps 2k and 4k.

\textbf{Interpretation.}  The two \(K\) differ in what the hard model does with its boundary.  At \(K=32\) of 4096 the
boundary sits in the tail of the score distribution, the competition for the last slots stays close, and
information about the candidates just below it (what the support term supplies) is what the hard gradient lacks.
At \(K=256\) the hard model separates the selected features from the rest: their values grow under the hard
gradient while the others get none, and the boundary region empties.  The support term works against that
separation --- its corrections are exchanges that pull both sides toward the boundary --- and a crowded boundary
in a hard forward means unstable selections and slots held by weak features (Entry 8 above).  The surrogate helps
where the selection is competitive and hurts where it can separate; at a given \(N\) which regime holds depends on
\(K\).

\textbf{Measured: the active-member term at \(K=256\).}  Restricted to the inactive members, the first-order estimator
at \(K=256\) is \(+0.041\), \(+0.036\), \(+0.026\), \(+0.020\) behind the hard backward at steps 500--2000, closing; its code
looks like the hard model's: 5 / 7 / 13 always-on features at gates 6 / 12 / 18 (hard 5 / 10 / 15, full pool 164 /
153 / 67) and a boundary at 3.0--3.5 (hard 2.7--3.2, full pool 1.6).  At \(K=32\) the same restriction loses the whole gain: \(+0.016\),
\(+0.021\), \(+0.007\), \(+0.001\) at steps 500--2000 (val CE 2.105 at step 2000), against \(-0.076\) for the full pool.  The scale-free kernel does not help at
\(K=256\) (\(+0.092\), \(+0.109\), \(+0.096\), \(+0.078\) at steps 500--2000).  At \(T=2\) the first-order estimator gains more at
\(K=32\): \(-0.088\) at step 2000 (\(T=1\): \(-0.076\)).

\textbf{Interpretation.}  The term on active features --- the one that pushes marginal selected features toward the
boundary --- is what crowds the boundary at \(K=256\): without it the code keeps the hard model's separation and the
deficit closes; with it the surrogate is best at \(K=32\), where the selection never separates.  The candidates'
term alone is the safe part of the surrogate; the active features' term pays only where the competition for the
last slots stays close.

\textbf{The \(K=256\) runs to 20k} (val CE, difference to the hard backward 1.3067): first order 1.3844 (\(+0.078\);
\(+0.069\) at 10k, \(+0.075\) at 16k); first order on the inactive members 1.2998 (\(-0.007\); \(+0.004\) at 10k, \(-0.000\)
at 12k, \(-0.006\) at 16k).  \textbf{Measured:} the full first-order estimator does not close its gap at \(K=256\); the
inactive-member version crosses near step 12k and ends slightly ahead.  At \(K=32\) the inactive-member version had
no gain at 2k; the variant that helps depends on \(K\).

% ============================================================
\section{Summary}
% ============================================================

\begin{enumerate}[leftmargin=1.4em]
\item \textbf{Why the unmodified surrogate fails.}  The code residual carries each code coordinate through an
  identity, and the unmodified backward applies every gate's support term to the total upstream, which already
  contains the terms of the gates above: along a feature carried near the boundary they multiply,
  \(\prod(1+|u|\kappa)\approx3^n\) (Entries 1--2).  These products are terms of second and higher order in the changes
  of support, which a hard forward does not have.  The gradient reaching the entry gate is \(2.6\cdot10^5\) times the
  hard model's at initialization; in training the code collapses onto always-on features and dead ones (Entry 6).
\item \textbf{What the existing knobs do.}  The support scale reduces the products to order \(\gamma^2\) and trains for
  \(\gamma\le0.3\) (\(\gamma=0.3\): \(-0.040\) at 2k, shrinking to \(-0.028\) at 8k; Entry 5).  The temperature, the window and
  the boundary mode leave the products in place (Entry 1); a scale-free kernel does not stop the routed collapse
  (Entry 5).
\item \textbf{A tempting fix that fails.}  Routing the support term into the update only removes the products but
  gives the carry and the update different gradients for one input; the code fills with always-on features
  (Entries 4, 6).  Centering the support term over tokens breaks the kernel's locality and fails faster (Entry 6).
\item \textbf{The principled fix.}  The first-order estimator (\code{rblapsum\_surrogate\_scope=first\_order}): the hard
  gradient plus, for each gate, the change its support term makes when every other gate is exact; one backward pass
  through each gate's input (R1) and no products (R2) (Entry 7).  Best at \(K=32\) (\(-0.076\) at 2k; \(-0.064\) at
  \(J=480\)); the advantage shrinks with training (\(-0.065\) at 5k).
\item \textbf{To 20k steps.}  At \(K=32\) the first-order estimator ends 0.045 below the hard backward (1.457 against
  1.502), the damped unmodified backward 0.016 below (scale 0.3).  At \(K=256\) the full first-order estimator stays
  0.078 behind; restricted to the inactive members it ends 0.007 ahead.  The term on active features helps where the
  hard model's selection stays competitive at the boundary (\(K=32\)) and hurts where the hard model separates its
  selected features (\(K=256\)) (Entry 8).
\end{enumerate}

\end{document}
